% CAISc 2026 submission — Open-Ended Track.
% Anonymized for review (default submission mode).
%
% Build: pdflatex paper && bibtex paper && pdflatex paper && pdflatex paper

\documentclass{article}

\PassOptionsToPackage{numbers,sort&compress}{natbib}
\usepackage{caisc_2026}

\usepackage[utf8]{inputenc}
\usepackage[T1]{fontenc}
\usepackage{hyperref}
\hypersetup{colorlinks=true, citecolor=blue, linkcolor=blue, urlcolor=blue}
\usepackage{url}
\usepackage{booktabs}
\usepackage{amsfonts}
\usepackage{amsmath}
\usepackage{amssymb}
\usepackage{nicefrac}
\usepackage{microtype}
\usepackage{xcolor}
\usepackage{graphicx}
\usepackage{algorithm}
\usepackage{algpseudocode}
\algrenewcommand\algorithmiccomment[1]{\hfill\textcolor{gray!60!black}{\(\triangleright\)\,\textit{#1}}}
\title{Verification-Gated Autonomous Optimization of Systems Code}

% Author block stays anonymized for submission (the .sty handles this
% automatically in the default mode). Do not list AI systems here; the
% AI Involvement Checklist at the end captures their role.
\author{%
  Anonymous Author(s) \\
  Affiliation \\
  Address \\
  \texttt{email}
}

\begin{document}

\maketitle

\begin{abstract}
We study whether an autonomous, LLM-driven optimization loop can
discover algorithmic improvements to a production consensus library
while preserving every safety property of the protocol. Our target is
\texttt{etcd-io/raft}, the consensus library beneath Kubernetes and
CockroachDB. We build a layered verification harness around it: an
existing TLA$^{+}$ specification for protocol-level invariants, a
deterministic fault-injection simulator that drives the library
through partition, drop, reorder, and duplication scenarios, and a
linearizability check on the resulting client histories. The harness
then gates an evolutionary search driven by a frontier coding model:
every candidate edit must clear compilation, the fault sweep, and
the linearizability check before its benchmark score is even
reported.
Our central claim is methodological: where safety is
mechanically checkable, the harness (not the model) is the
load-bearing artefact, and a sufficiently sharp harness lets a
general-purpose coding agent search inside a mature systems codebase
safely. Harness-acceptance is necessary but not sufficient for
upstream merge; we make no production-readiness claim.
\end{abstract}

\section{Introduction}
\label{sec:intro}

Raft \cite{ongaro2014search} is the consensus algorithm at the core of
most contemporary distributed systems: etcd~\cite{etcd},
Kubernetes~\cite{kubernetes}, CockroachDB~\cite{cockroachdb},
Docker Swarm~\cite{dockerSwarm}, TiKV~\cite{tikv}. Implementations of
Raft are required to preserve
five safety properties (election safety, log matching, leader
completeness, state-machine safety, and log durability under partial
failure) under every interleaving of message loss, reorder, duplication,
and crash-restart. The cost of an undetected regression is, in
practice, total: a single safety-violating commit can corrupt the entire
replicated state machine across operators.

Recent work on AI-driven research for systems
\cite{barbarians2025,datadogHarness,datadogAutonomous} has shown that
LLM-driven evolutionary search can match or exceed human-tuned
performance on isolated kernels, schedulers, and codec libraries when
correctness is verifiable. The natural question is whether the same
loop can be turned on a production-grade consensus library, where
correctness is a global protocol invariant rather than a per-operation
predicate.

Our hypothesis is that the bottleneck is verification bandwidth, not
generation quality. Concretely: if every candidate variant produced by
the LLM is forced through a mechanical pipeline (compile, seeded
deterministic fault-injection sweep, linearizability check,
benchmark) before any performance score is reported, the
generator can be allowed to explore aggressively without endangering
the codebase. The harness, not the model, becomes the artefact whose
quality determines what kinds of improvements are reachable.

This paper instantiates that pattern over
\texttt{etcd-io/raft}~\cite{etcdraft} at a pinned commit. Our contributions are:
\begin{enumerate}
  \item A minimal-upstream-patch deterministic simulation testing
    (DST)~\cite{fdbsim,foundationdb} harness over the library's existing
    \texttt{RawNode}/\texttt{Storage} interfaces, with
    seeded fault injection (partition, heal, drop, reorder,
    duplicate) and a three-layer oracle (TLA$^{+}$ trace validation, Raft-protocol invariants and
    Porcupine linearizability on client histories). The harness
    achieves bit-for-bit reproducibility per seed after a four-line
    patch to the upstream election-timeout random number generator
    (RNG).
  \item Calibration of harness noise floors: we report 0~invariant
    violations across 200 distinct fault-injected seeds on the
    unmodified upstream, a 17\% Porcupine false-positive rate
    attributable to tick-level read-resolution timing, and a
    bimodal benchmark distribution.
  \item An evolutionary-search wrapper that uses the harness as the
    rejection oracle. A frontier coding model proposes bounded edits
    to the leader-side flow-control function
    \texttt{maybeSendAppend}; candidates that fail to
    compile, fail any upstream test, fail any fault-injection
    scenario, or produce non-linearizable client histories beyond
    the noise floor are discarded before benchmarking.
  \item An empirical study reporting accepted variants, their
    throughput characteristics against the pinned baseline, and an
    ablation indicating which harness layer first catches which
    class of incorrect candidate.
\end{enumerate}

Non-goals.  We do not retrain or fine-tune the underlying model.  We
do not propose accepted variants as production-ready: harness
acceptance is necessary but not sufficient for upstream merge, and
human review and shadow deployment remain in scope.

\section{Background}
\label{sec:background}

\subsection{Raft}
Raft \cite{ongaro2014search} is a leader-based consensus protocol over
a cluster of replicas. A unique leader per term linearizes all
proposals via an append-only replicated log; a majority of replicas
must persist an entry before it is committed; followers apply
committed entries to an arbitrary state machine in log order. The
protocol's safety rests on five invariants: \emph{election safety}
(at most one leader per term); \emph{leader append-only}; \emph{log
matching} (identical prefixes across replicas); \emph{leader
completeness} (any committed entry appears in every subsequent leader
log); and \emph{state-machine safety} (no two replicas apply
divergent commands at the same index). We treat all five as required
invariants in our harness.

\subsection{The \texttt{etcd-io/raft} Library}
\label{sec:bg-etcd}
\texttt{etcd-io/raft} is a Go library implementing Raft as a pure
state machine: it owns no network sockets, no disk, and no goroutine
loops. The public \texttt{Node}/\texttt{RawNode} surface exposes
\texttt{Tick()}, \texttt{Step(m)}, \texttt{Propose}, \texttt{ReadIndex},
\texttt{Campaign}, \texttt{Ready()}, and \texttt{Advance}; a pluggable
\texttt{Storage} interface provides log persistence. This design makes
the library unusually amenable to deterministic simulation testing:
ticks are logical, message delivery and disk are caller-controlled,
and the state machine's output is fully determined by its input
sequence, given that the few sources of non-determinism inside the
library are themselves controllable.

The repository ships three artefacts we reuse. (1)~A \texttt{rafttest}
package providing a command-driven interaction environment for
data-driven tests; we adopt its \texttt{InteractionEnv} only as a
reference, building our own scheduler atop \texttt{RawNode} for tighter
control. (2)~A \texttt{state\_trace.go} file under a
\texttt{with\_tla} build tag emitting \texttt{TracingEvent} records at
every state transition. (3)~A \texttt{tla/} directory containing
\texttt{MCetcdraft.tla} (the model-check spec) and
\texttt{Traceetcdraft.tla} (a trace-validation spec consuming the
NDJSON output of \texttt{state\_trace.go}).

\subsection{Harness-First and Verification-Gated Search}
The harness-first thesis~\cite{datadogHarness} inverts the
model-vs-review tradeoff: instead of asking whether agent-written
code is correct, a layered set of automated oracles asks whether the
candidate falsifies any property. Prior instantiations target a
Redis-compatible store~\cite{redisRust} and a Kafka-compatible
engine~\cite{datadogHarness}. Verification-gated autonomous
optimization~\cite{datadogAutonomous,barbarians2025} extends the
pattern with an LLM candidate generator and an evolutionary
database~\cite{openevolve,alphaevolve,funsearch,alphadev}, gated on a
formal or empirical oracle; reported deployments have driven large
speedups on time-series codecs behind \texttt{Verus}~\cite{verus}
safety proofs. We instantiate the pattern over a Go
consensus library: the harness combines TLA$^{+}$ trace validation,
deterministic simulation, and linearizability checking, and the
search target is a single function inside a $\sim$85~KB host file.

\section{Verification Harness}
\label{sec:harness}

\subsection{Design Principles}
We adopt three principles. \emph{Independence}: each layer should
fail for a different reason: a wrong protocol spec passes
deterministic simulation, an incomplete fault catalogue passes the
spec, a permissive state-machine model accepts non-linearizable
histories that the spec rejects. \emph{Reuse-first}: where upstream
already ships verification infrastructure (the \texttt{state\_trace.go}
hooks and the in-tree TLA$^{+}$ files), we wrap rather than
reimplement. \emph{Cheap iteration}: every layer must terminate in
seconds per scenario, since the optimization loop in \S\ref{sec:loop}
will invoke the full pyramid hundreds of times per run.

We achieve bit-for-bit reproducibility per seed via a single
four-line patch to the upstream library: a new optional
\texttt{Config.Rand} field that, when set, replaces a
\texttt{crypto/rand}-backed source for the randomized election
timeout (the only non-seedable randomness in the library, identified
by a single \texttt{globalRand.Intn} call in \texttt{raft.go}). When
the field is nil, the existing behaviour is preserved; the upstream
test suite continues to pass without modification.

\subsection{Layer 1: Protocol Specification}
\label{sec:harness-tla}
For the protocol-level oracle, we reuse the in-tree
\texttt{MCetcdraft.tla} spec and the \texttt{Traceetcdraft.tla}
trace-validation spec. The harness is built with
\texttt{-tags=with\_tla} so that \texttt{state\_trace.go} emits a
\texttt{TracingEvent} stream on every state transition (election,
becoming leader, sending append entries, etc.). A run's NDJSON trace
is then validated against the TLC model checker via the upstream
\texttt{tla/validate.sh} script; any state or transition not
admitted by the spec is a failure.

\subsection{Layer 2: Deterministic Simulation Testing}
\label{sec:harness-dst}
The DST layer is a from-scratch single-goroutine scheduler built atop
\texttt{RawNode}; it advances the cluster one logical tick at a time
in four phases.
\textbf{Clocks:} every healthy node is ticked, possibly firing an
election or a heartbeat.
\textbf{Progress:} each node drains its \texttt{Ready()}; the
harness persists the new entries to a per-node \texttt{MemoryStorage},
applies committed entries to a single-register state machine, and
enqueues outbound messages onto the network in a canonical order.
A deterministic sort is required because upstream builds the outbox
by iterating a Go map.
\textbf{Delivery:} messages whose scheduled-delivery tick has arrived
are \texttt{Step()}-ed into their destination, possibly being
dropped, reordered, or duplicated by the network layer.
\textbf{Reads:} pending client reads complete once the linearizable
safe index has caught up.
All randomness (ours and upstream's election timeout) flows
from a single seeded \texttt{*rand.Rand}, exposed via the four-line
patch of \S\ref{sec:harness}; full pseudocode and the bindings to
upstream's \texttt{RawNode}/\texttt{Storage} interface are in
Appendix~\ref{app:dst-detail}.

The network supports per-link drop probability, partition/heal,
per-message scheduled-delivery delay (yielding reorder when delays
exceed one tick), and per-send duplication. The default scenario
applies $p_\text{drop}=0.02$ on every directed link,
$p_\text{partition}=0.01$ per tick, $p_\text{heal}=0.05$,
$\text{maxDelay}=3$~ticks, and $p_\text{dup}=0.02$.

The invariant oracle checks two properties on every tick:
\emph{log-prefix agreement} (for any two nodes $i,j$ and any index $k$
present in both node's committed list, the entries must be identical)
and \emph{leader uniqueness per term} (no two nodes are in
\texttt{StateLeader} at the same term).

\subsection{Layer 3: Linearizability}
\label{sec:harness-lin}
We adopt the classical linearizability
correctness condition~\cite{herlihy1990linearizability} on the
sequence of completed client operations. Three logical clients, each
preferring a distinct cluster node, issue writes via \texttt{Propose}
and linearizable reads via \texttt{ReadIndex} into a single-register
state machine. Each client
keeps at most one operation in flight; writes encode a 17-byte payload
 so committed entries can be matched
back to the issuing client. Reads are finalised once the preferred
node's \texttt{appliedIndex} crosses the safe index returned in
\texttt{Ready.ReadStates}; the read returns the register snapshot
captured at the safe index (rather than the current register), which
is necessary because multiple entries may apply in a single
\texttt{Ready} cycle. The recorded $\langle$invoke, return$\rangle$
history is then checked against a register specification using
\texttt{Porcupine}~\cite{porcupine}, a Go re-implementation of the
checker pattern introduced by \texttt{Knossos}~\cite{knossos} for
Jepsen~\cite{jepsen}.

\subsection{Fault Catalogue and Coverage}
The catalogue is intentionally small but distinguishing: partition,
heal, per-link drop, per-message delay/reorder, and per-send
duplication. Crash-restart is excluded by construction because our
\texttt{MemoryStorage} does not persist conf state without snapshot
support; we treat this as a deferred extension and discuss it in
\S\ref{sec:limits}.

\subsection{Layer Costs}
\label{sec:harness-cost}
On a 16-vCPU \texttt{x86\_64} cloud host ($32$~GB RAM, Ubuntu
\texttt{24.04}), single-process, the per-candidate cost of the
harness pyramid (compile, $8$ fault-injected DST seeds at $1500$~ticks
each, Porcupine check, and a five-run benchmark) is approximately
$15$--$20$~seconds wall-clock; the LLM call dominates total
per-iteration time. The DST layer alone runs at $\sim$12\,k
ticks/second for a three-node cluster; a 200-seed sweep at
$2000$~ticks completes in under four minutes.

\section{Autonomous Optimization Loop}
\label{sec:loop}

\subsection{Search Space}
We restrict the optimizer to one function in \texttt{raft.go}, the
leader's per-follower send decision
\texttt{maybeSendAppend(to, sendIfEmpty)}. It runs once per follower
per tick and decides whether to emit an \texttt{AppendEntries}
message, governing flow control, batching, pipelining, and snapshot
fallback, the hottest function on the leader's outbound path
($\sim$45~lines). We chose it because every throughput lever is local
to the function, yet a wrong edit can break log consistency, making it
a sharp test of the harness rather than the generator. The levers in
scope (a throttle check, the replication-state predicate, the
owed-entries log slice, the post-send progress update, and the
snapshot fallback) map to their upstream identifiers in
Table~\ref{tab:levers}. The mutable region is delimited by
\texttt{EVOLVE-BLOCK} markers; diffs touching the file outside it are
rejected before compilation.

The generator is instructed to preserve five behavioural contracts,
stated by purpose: the signature; the throttle short-circuit (never
send to a paused or window-full follower); snapshot fallback when the
required log prefix has been compacted; the post-send progress update
(else later ticks break log matching); and no new package-level state,
goroutines, locks, or I/O. The middle three are exactly what the
harness re-checks mechanically: the prompt states them so the
generator need not rediscover them, but the harness, not the prompt,
enforces them. The verbatim prompt is in
Appendix~\ref{app:system-prompt}.

\subsection{Candidate Generator}
A frontier coding agent receives the entire host file plus the
system prompt above and is asked to emit a SEARCH/REPLACE diff scoped
to the marked region. The generator is invoked through an
OpenAI-compatible chat endpoint and is, by configuration, never
informed of the benchmark scores produced by previous candidates;
only that they were ``accepted'' or ``rejected'' by the harness, with
diagnostics from the failing layer if any. This design keeps the
generator from over-fitting to bench noise.

\subsection{Rejection Oracle}
The rejection oracle is the harness pyramid of \S\ref{sec:harness},
applied in order of increasing cost: \emph{(a)}~upstream package
compilation; \emph{(b)}~the DST sweep ($N=8$ seeds, $1500$~ticks
each, full fault catalogue); \emph{(c)}~the linearizability check on
the recorded client histories; only then \emph{(d)}~the benchmark.
Any layer's failure terminates evaluation immediately and is
recorded as the candidate's rejection reason. The harness layers
contribute failure modes independently: a candidate that compiles
correctly but violates log-matching is caught at (b), while a
candidate that passes (b) but reorders \texttt{Ready.Messages} in a
way that breaks linearizability is caught at (c). The unmodified
upstream test suite is \emph{not} run per candidate (it adds
$\sim$5~seconds for negligible additional coverage given (b)) but
is run on accepted candidates as part of the post-acceptance
re-validation described in \S\ref{sec:results-variants}.

\subsection{Candidate Database and Selection}
We use a published open-source evolutionary search
framework~\cite{openevolve} that maintains a SQLite-backed candidate
database with MAP-Elites-style island distribution. Parent sampling
is power-law over accepted candidates; mutation is the LLM's
SEARCH/REPLACE diff. The framework retains rejected candidates with
their rejection-reason metadata, which is consumed by the per-layer
ablation in \S\ref{sec:results}.

\subsection{Scoring}
A candidate that clears the oracle is scored as
\[
  s = \frac{n^{\text{base}}_\text{op}}{\min_{i=1\ldots K}\, n^{\text{cand},i}_\text{op}}
\]
where $n^{\text{cand},i}_\text{op}$ is the per-run ns/op of the
upstream \texttt{BenchmarkProposal3Nodes} on the candidate, $K=5$
runs are taken to minimise jitter sensitivity, and
$n^{\text{base}}_\text{op}$ is a fixed calibration constant chosen so
the unmodified baseline scores near $1.0$. Scores below the
calibrated noise floor are not interpretable as improvements; see
\S\ref{sec:results} for the floor measurement.

\section{Experimental Setup}
\label{sec:setup}

\paragraph{Target.}
\texttt{etcd-io/raft} pinned at commit \texttt{67d129e8} (head of
\texttt{main} on 2026-05-26). The single patch to upstream is the
seedable election-timeout RNG of \S\ref{sec:harness}; all other
modifications live in a separate Go module and are applied to a
working copy of \texttt{raft.go} on each evaluation. The unmodified
upstream test suite is run on accepted candidates as part of
post-acceptance re-validation.

\paragraph{Software and hardware.}
Go \texttt{1.26.3}; OpenEvolve~\cite{openevolve} \texttt{0.2.27};
Porcupine~\cite{porcupine} \texttt{v1.1.0}; OpenJDK \texttt{17} with
\texttt{tla2tools}. The headline benchmark
(Table~\ref{tab:bench-paired}) ran on a dedicated 16-vCPU
\texttt{x86\_64} cloud host with $32$~GB RAM, Ubuntu \texttt{24.04}
LTS, no co-tenant workload. All harness sources, the patch,
candidate-database snapshots, prompt templates, and the runbook are
released as the verification artefact.

\paragraph{Workload.}
Three logical clients, one preferring each node of a three-node
cluster, issue mixed reads (30\%) and writes (70\%) with per-client
$p_\text{issue}=0.3$ per tick, gated on a 200-tick per-operation
timeout. Per-tick fault injection probabilities and the message
delay/duplication parameters are as in \S\ref{sec:harness-dst}.

\paragraph{Seeds.}
The DST sweep uses a deterministic set of seeds
$\{1,\ldots,N\}$ with $N=8$ per candidate evaluation. For the
calibration measurements in \S\ref{sec:harness} we use $N\in\{50, 100,
200\}$ at $2000$~ticks.

\paragraph{Benchmark.}
We invoke \texttt{go test -bench=BenchmarkProposal3Nodes -benchtime=1s
-count=5} and report the minimum ns/op across the five runs; the
minimum is more robust than mean or median against background
system-load excursions.

\paragraph{LLM budget.}
The evolution loop is bounded by a per-run iteration count and a
per-iteration wall-clock cap. The frontier coding model is accessed
through an OpenAI-compatible chat endpoint; we report token spend
and total wall-clock alongside results in \S\ref{sec:results}.

\section{Results}
\label{sec:results}

\subsection{Calibration and Noise Floors}
On the unmodified upstream library at the pinned commit, the invariant
oracle of \S\ref{sec:harness-dst} reports $0$~violations across
$200$~seeds at $2000$~ticks each with the full fault catalogue
enabled. With faults disabled, the linearizability check passes
$29$ of $30$ seeds ($97\%$); with faults enabled, the pass rate falls
to $83$ of $100$ seeds. The $17\%$ false-positive rate under faults
is attributable to tick-level read-resolution timing (multiple
entries applying in one \texttt{Ready} cycle with overlapping
write/read intervals) and leader-change races during in-flight reads;
the harness documents this floor and uses it as the threshold below
which linearizability-check failures are not treated as candidate
rejections in the loop. We treat the invariant oracle as the strict
gate.

The benchmark exhibits a strongly bimodal ns/op distribution under
repeated runs of an unmodified binary. A $100$-sample calibration on
the evaluation host (Appendix~\ref{app:calibration}) gives a
fast-mode minimum of $4905$~ns/op against a median of $8972$ and a
p90 of $10734$, a spread driven by system load; the in-loop evaluator
reports the minimum across five \texttt{-benchtime=1s} runs, a choice
we revisit in \S\ref{sec:results-variants}. The in-loop score divides
by a per-host baseline from a coarser earlier calibration ($2400$~ns/op
arm64, $4688$~ns/op x86); $4905$ is a later $100$-sample recalibration,
and as a common divisor the baseline does not affect rankings.

\subsection{Loop Output}
\label{sec:results-loop}
We ran the loop for $30$ iterations over $\sim$28 minutes of wall-clock
on a single commodity arm64 host (in-loop baseline $2400$~ns/op). Of
$30$ proposed candidates, \emph{all $30$ cleared
the rejection oracle}: every candidate compiled, passed $8$
fault-injected DST seeds at $1500$~ticks, and produced a
non-violating client history. None were
rejected at the linearizability layer (the gate operates at the
calibrated noise floor of \S\ref{sec:harness-lin}). The mean per-iteration
score was $1.067$; $25$ of $30$ candidates ($83.3\%$) scored at or
above the calibrated baseline, with $5$ candidates regressing
($\min=0.795$). The best score observed during the loop was
$1.1516$ at iteration~10 (program identifier
\texttt{1c3a355f}; bench evaluator $n_\text{op}=2084$~ns under
min-of-five sampling).

The strict invariant gate fired zero times across all $30$
candidates' fault-injection sweeps. This is not a claim of generator
infallibility; it is a claim about the LLM's tendency, under our
prompt scoping, to produce edits inside a code region where most
plausible mutations preserve the function's contract.

\paragraph{Reproducibility.}
An independent second $30$-iteration run on the x86 droplet (fresh
sampling, identical configuration and gate) reproduced the gate
behaviour exactly: all
$30$ proposed candidates again cleared compilation, the $8$-seed
fault-injected DST sweep, and the linearizability layer, with zero
invariant violations. We use this replicate for the scoring analysis
in \S\ref{sec:results-variants} and the ablation in
\S\ref{sec:results-ablation}.

\subsection{Best Accepted Variant}
\label{sec:results-variants}
The best accepted variant (ranked by paired median) introduces three
coordinated edits to \texttt{maybeSendAppend}:

\begin{enumerate}
  \item \textbf{Throttled early-bail}: the function computes
    \texttt{throttled := pr.State == StateReplicate \&\&
    pr.Inflights.Full()} once at entry; when
    \texttt{throttled \&\& !sendIfEmpty} the function returns
    immediately, skipping the storage-backed
    \texttt{raftLog.term()} lookup and all subsequent log access.
  \item \textbf{Skip \texttt{payloadsSize} on empty entries}: the
    common heartbeat-style empty-MsgApp path bypasses
    \texttt{payloadsSize([])} entirely.
  \item \textbf{Local cleanup}: \texttt{pr.Next} and
    \texttt{r.raftLog.committed} are hoisted into stack-locals, and
    the now-dead post-loop \texttt{if err != nil} branch is
    removed.
\end{enumerate}

We re-validated the variant against three independent oracles. The
full upstream test suite passes on all six packages
(\texttt{go test ./...}). A $50$-seed harness sweep at $2000$~ticks
each, with the full fault catalogue, reports $0$ invariant
violations. Finally, a paired $10$-run benchmark using
\texttt{-benchtime=2s} on a dedicated commodity cloud host between
back-to-back runs of \texttt{raft.go.baseline} and the candidate
yields the distribution in Table~\ref{tab:bench-paired}. The
candidate is substantially faster across the body of the
distribution: \emph{$25.4\%$ median improvement
($9405 \rightarrow 7020$~ns/op), $26.1\%$ at p25, $17.4\%$ at the
max}. The candidate is slightly slower at the very minimum
($-7.8\%$), reflecting a small fixed overhead from the extra
throttle check; it more than compensates by collapsing the
distribution's long tail.

\begin{table}
  \caption{Paired $10$-run \texttt{BenchmarkProposal3Nodes} on a
    dedicated 16-core Intel host: ns/op (lower is better).}
  \label{tab:bench-paired}
  \centering
  \begin{tabular}{lrrrr}
    \toprule
    & min & p25 & median & max \\
    \midrule
    baseline   & $5683$ & $8678$ & $9405$ & $11439$ \\
    candidate  & $6127$ & $6416$ & $7020$ & $9448$ \\
    \midrule
    improvement & $-7.8\%$ & $+26.1\%$ & $+25.4\%$ & $+17.4\%$ \\
    \bottomrule
  \end{tabular}
\end{table}

\paragraph{Evaluator-side score under-ranked the headline candidate.}
The per-iteration evaluator scored this candidate at
$\approx\!0.998$ (a marginal regression) because its scoring
function takes the $\min$ of five back-to-back bench runs. The bench
has a wide bimodal distribution on this host; $\min$-of-$N$ is
dominated by lucky-fast runs, where the candidate's small fixed
overhead shows. The paired-distribution analysis in
Table~\ref{tab:bench-paired} reveals that the candidate's
optimisations affect the tail far more than the floor, exactly the
behaviour one wants in a consensus library, where p99 latency, not
best-case throughput, governs end-to-end commit time.

The independent replicate of \S\ref{sec:results-loop} makes the
effect quantitative. Re-scoring its $31$ evaluated candidates against
the same calibration baseline, $\min$-of-$N$ counts only $3$ as
improvements while the paired median counts $11$; more starkly, the
two rules invert the ranking. The loop's $\min$-best candidate
(\texttt{90f41d5d}, $\min$-speedup $1.06$) is a median \emph{regression}
($0.98\times$, ranked $18$th of $31$ by median), whereas the genuine
median-best (\texttt{5866c047}, $1.60\times$ on the median by
collapsing the bimodal tail) is ranked only $9$th by $\min$ and would
be discarded as a regression ($\min$-speedup $0.92$). The full
re-ranking is in Appendix~\ref{app:reranking}. A more robust evaluator
would therefore score on the paired median rather than the min; we
treat this as a methodological lesson and discuss it in
\S\ref{sec:limits}.

\subsection{Per-Layer Ablation}
\label{sec:results-ablation}
Across both $30$-iteration runs the strict gate never fired, showing
the constrained prompt is safe but not that the four layers are
\emph{non-redundant}. To measure layer-uniqueness we hand-wrote five
faulty variants of \texttt{maybeSendAppend}, each breaking a single
named contract, plus an unmodified control, and pushed every variant
through all four layers at $N=12$ seeds. The full outcome is in
Appendix~\ref{app:ablation} (Table~\ref{tab:ablation-full}); it has
three parts.

First, \emph{every} broken variant is rejected before it can be
benchmarked, and the catches span all four layers: a
nonexistent-method call at compile; an omitted progress update by the
DST sweep (a seed-1 liveness timeout); a forged commit index by the
linearizability checker; and a paused-follower flood and a
\texttt{prevIndex} off-by-one by the upstream suite (the flood also
tripped linearizability at $2/12$ seeds, at the $1/12$ noise floor, so
L3 is what catches it reliably). Second, the
layers are genuinely independent: the safety-critical commit-index
lie is flagged non-linearizable on $10/12$ seeds (far above the
control's $1/12$ noise floor), yet \emph{passes} the invariant oracle
untouched, precisely the redundancy Layer~3 exists to provide. Third,
the harness does not subsume the upstream suite: a self-correcting
\texttt{prevIndex} off-by-one passes both DST layers untouched and the
paused-follower flood clears them up to the noise floor, yet both are
caught reliably by upstream unit tests, which is why we re-run that
suite on accepted candidates (\S\ref{sec:results-variants}).

\subsection{Cost}
\label{sec:results-cost}
The $30$-iteration loop completed in approximately $28$~minutes
wall-clock, dominated by per-iteration evaluation cost
(median $\sim$57~s; $\sim$15~s LLM call, $\sim$3~s upstream compile,
$\sim$8~s harness compile, $\sim$1~s DST sweep at $N=8$ seeds and
$1500$~ticks, $\sim$30~s benchmark with five runs). Each iteration
sends the full host file (the $\sim$85~K-byte \texttt{raft.go} with
markers) plus the system prompt as input to the LLM and receives a
SEARCH/REPLACE diff back. Token spend per iteration is on the order
of $2$\,$\cdot 10^{4}$ input tokens and $5$\,$\cdot 10^{2}$ output tokens,
yielding a total spend bounded by approximately ten US dollars at
publicly listed inference prices for the frontier model class used.
The harness, not the model, dominates wall-clock; halving the
benchmark time would roughly halve the loop's run time without
materially affecting the result.

\section{Limitations and Threats to Validity}
\label{sec:limits}

\paragraph{Oracle coverage is a lower bound.}
The invariant oracle covers log-prefix agreement and per-term leader
uniqueness; it does not cover lease-based read safety,
snapshot-and-truncation correctness, or membership-change conflict
avoidance. The linearizability check uses a single-register state
machine. The fault catalogue covers partition, drop, reorder, and
duplication; crash-restart and slow-disk are scheduled extensions.

\paragraph{Synthetic performance, harness-only acceptance.}
All benchmarks are against an in-memory simulated network and
\texttt{MemoryStorage}; we make no claim about production traffic or
physical contention. Harness-acceptance is necessary but not
sufficient for upstream merge; mainline review and shadow deployment
remain in scope.

\paragraph{Evaluator scoring and generator noise.}
$\min$-of-$N$ benchmark scoring during evolution is jitter-sensitive
and demonstrably misranked our best candidate; paired-median scoring
(Table~\ref{tab:bench-paired}) is more robust and should be the in-loop
metric in future runs. The coding model is stochastic and closed-source;
different prompts, temperatures, or checkpoints will explore different
regions.

\section{Related Work}
\label{sec:related}

\paragraph{Harness-first and verification-gated.}
The methodological antecedents are the harness-first
\texttt{redis-rust}~\cite{redisRust} and Helix~\cite{datadogHarness}
studies, and the Unicron autonomous-optimization
pipeline~\cite{datadogAutonomous} which gates LLM-evolved Rust
behind a \texttt{Verus}~\cite{verus} safety proof. Our work
specialises the pattern to a single function in a production Go
library (no formal verification gate; an empirical DST oracle in its
place) and makes the harness the contribution rather than the
implementation.

\paragraph{Formal verification of distributed systems.}
\texttt{IronFleet}~\cite{ironfleet} (Dafny) and
\texttt{Verdi}~\cite{verdi} (Coq) gave machine-checked proofs of
distributed implementations at scale, and \texttt{Verus}~\cite{verus}
brings SMT-backed verification to Rust. Go has no such analogue, so
our Layer~1 is an existing TLA$^{+}$ trace validator rather than a
refinement proof.

\paragraph{LLM-driven code synthesis and evolution.}
The lineage runs from \texttt{Codex}~\cite{codex} and
\texttt{AlphaCode}~\cite{alphacode} to evolutionary search
(\texttt{FunSearch}~\cite{funsearch},
\texttt{AlphaDev}~\cite{alphadev},
\texttt{AlphaEvolve}~\cite{alphaevolve}).  The ADRS
programme~\cite{barbarians2025} applies the latter via an open-source
framework~\cite{openevolve} across ten benchmarks; we narrow it to one
safety-critical function. \texttt{SWE-bench}~\cite{swebench}-style agent
benchmarks gate at the repository level, not the function-and-trace
level we use.

\paragraph{Deterministic simulation and distributed-systems testing.}
FoundationDB's simulator~\cite{fdbsim,foundationdb} and
Jepsen~\cite{jepsen} pioneered seeded simulation and linearizability
checking, continued by \texttt{Antithesis}~\cite{antithesis},
\texttt{TigerBeetle}'s \texttt{VOPR}~\cite{tigerbeetleDST}, and
\texttt{Madsim}~\cite{madsim}, atop the property-based seeded-random
discipline~\cite{quickcheck,hypothesis}. Our DST inherits FDB's
single-goroutine seeded-RNG style and the register-history checker of
Knossos~\cite{knossos}/Porcupine~\cite{porcupine}; \texttt{etcd}'s
robustness suite~\cite{etcdRobustness} targets the integrated server,
ours the algorithm library.

\section{Conclusion}
\label{sec:conclusion}

In a domain where safety properties are mechanically checkable, the
harness is the load-bearing artefact. Our pyramid gates an
LLM-evolutionary loop over one function in \texttt{etcd-io/raft};
the strict gate fires zero false positives across $200$~seeds, and
the headline accepted candidate delivers a $25.4\%$ paired-median
throughput improvement. The methodology generalises to any systems
library with a pluggable side-effect interface and a TLA$^{+}$
spec.

\begin{ack}
% Acknowledgements are hidden in the anonymized submission via the
% style file's `ack` environment; they appear only in the camera-ready
% copy.
\end{ack}

\section*{References}
\medskip
{
\renewcommand{\section}[2]{}% suppress thebibliography's auto-heading
\small
\bibliographystyle{plainnat}
\bibliography{references}
}

%%%%%%%%%%%%%%%%%%%%%%%%%%%%%%%%%%%%%%%%%%%%%%%%%%%%%%%%%%%%
\appendix
\section{Technical Appendices and Supplementary Material}

\subsection{DST scheduler: full algorithm and bindings to upstream
            \texttt{RawNode}}
\label{app:dst-detail}

Algorithm~\ref{alg:dst-tick} gives the per-tick step in full; the
remainder of this section explains how each line maps onto the
upstream \texttt{RawNode}/\texttt{Storage} interface.

\begin{algorithm}
\caption{One simulator tick (deterministic, single-threaded).}
\label{alg:dst-tick}
\begin{algorithmic}[1]
\Require cluster $C$, network $N$, seeded RNG $\rho$
\Statex \textit{1.\ Clocks.}
\ForAll{healthy $n \in C$}
  \State $n.\textsc{Tick}(\rho)$ \Comment{may fire election or heartbeat}
\EndFor
\Statex \textit{2.\ Let each node make progress.}
\ForAll{$n \in C$}
  \State $\langle\text{entries},\,\text{committed},\,\text{outbox}\rangle
         \gets n.\textsc{Ready}()$
  \State \textsc{Persist}(entries);\ \ \textsc{Apply}(committed);\ \
         enqueue \textsc{Sort}(outbox) into $N$
  \State $n.\textsc{Advance}()$
\EndFor
\Statex \textit{3.\ Deliver mail.}
\ForAll{$m \in N$ scheduled for this tick}
  \State $m.\text{dst}.\textsc{Step}(m)$
         \Comment{may drop, reorder, or duplicate}
\EndFor
\Statex \textit{4.\ Answer clients.}
\ForAll{pending read $q$ with $q.\text{idx} \le \text{safeIndex}$}
  \State complete $q$
\EndFor
\end{algorithmic}
\end{algorithm}

\paragraph{Phase 2 (Progress) bindings.}
The destructure
$\langle\text{entries},\text{committed},\text{outbox}\rangle$
corresponds to the \texttt{rd.Entries}, \texttt{rd.CommittedEntries},
and \texttt{rd.Messages} fields of the \texttt{Ready} struct returned
by \texttt{RawNode.Ready()}. \textsc{Persist} writes to a per-node
\texttt{MemoryStorage}; \textsc{Apply} updates a single-register
state machine; \textsc{Sort} orders the outbox by
$\langle\text{To},\text{From},\text{Type},\text{Term},\text{Index}\rangle$
before any message is enqueued. The sort is mandatory: upstream
constructs \texttt{r.msgs} by iterating a
\texttt{map[uint64]*Progress}, and Go intentionally randomises map
iteration order, so any deterministic execution must impose a
canonical order at this seam. After Phase 2 the harness calls
\texttt{n.Advance(rd)} to release the \texttt{Ready} block, matching
upstream's contract.

\paragraph{Phase 3 (Delivery) and randomness.}
The network applies drop, reorder, and duplication faults using the
parameters of \S\ref{sec:harness-dst}. All randomness (fault
choices, partition/heal events, scheduled-delivery delays, and
upstream's own election timeout) is sourced from a single seeded
\texttt{*rand.Rand}; the latter is exposed by the four-line patch of
\S\ref{sec:harness}, which is the only modification to the upstream
library.

\subsection{Search-Space Identifiers}
\label{app:search-identifiers}
Table~\ref{tab:levers} maps each lever named in \S\ref{sec:loop} to its
upstream \texttt{etcd-io/raft} identifier at the pinned commit.

\begin{table}[h]
  \caption{The optimization levers in \texttt{maybeSendAppend} and
    their upstream symbols.}
  \label{tab:levers}
  \centering
  \begin{tabular}{ll}
    \toprule
    Lever (\S\ref{sec:loop}) & Upstream identifier \\
    \midrule
    Throttle check     & \texttt{pr.IsPaused()}, \texttt{pr.Inflights.Full()} \\
    Replication state  & \texttt{pr.State == StateReplicate} \\
    Owed-entries slice & \texttt{r.raftLog.entries(pr.Next, r.maxMsgSize)} \\
    Progress update    & \texttt{pr.SentEntries(\ldots)}, \texttt{pr.SentCommit(\ldots)} \\
    Snapshot fallback  & \texttt{r.raftLog.term(pr.Next-1)} error path \\
    \bottomrule
  \end{tabular}
\end{table}

\subsection{Host Benchmark Calibration}
\label{app:calibration}
Table~\ref{tab:calibration} is the unmodified-baseline ns/op
distribution on the evaluation host (16-vCPU \texttt{x86\_64},
$100$~samples: $20$ cycles of \texttt{-count=5}). The gap between the
minimum and the median is the bimodality that the in-loop
$\min$-of-$N$ scoring exploits (\S\ref{sec:results-variants}).

\begin{table}[h]
  \caption{Baseline \texttt{BenchmarkProposal3Nodes} distribution
    ($100$ samples, ns/op; lower is better).}
  \label{tab:calibration}
  \centering
  \begin{tabular}{rrrrrrr}
    \toprule
    min & p10 & p25 & median & p75 & p90 & max \\
    \midrule
    $4905$ & $5425$ & $6502$ & $8972$ & $9329$ & $10734$ & $15153$ \\
    \bottomrule
  \end{tabular}
\end{table}

\subsection{Planted-Bug Ablation: Full Results}
\label{app:ablation}
Table~\ref{tab:ablation-full} expands \S\ref{sec:results-ablation}.
``inv'' is the invariant DST sweep (seed of first violation, or
\texttt{--} if none); ``lin'' is the Porcupine fail-count over $12$
seeds; ``tests'' is the upstream root-package suite. The first layer
to reject each variant is in the final column.

\begin{table}[h]
  \caption{Planted-bug ablation, full per-layer outcomes ($N=12$
    seeds). The control is the unmodified baseline.}
  \label{tab:ablation-full}
  \centering
  \small
  \begin{tabular}{llccll}
    \toprule
    Variant (broken contract) & compile & inv & lin & tests & First catch \\
    \midrule
    control (none)                 & ok   & --    & $1/12$  & pass & (passes) \\
    nonexistent method call        & FAIL & --    & --      & --   & L0 compile \\
    omit progress update           & ok   & seed~1$^\dagger$ & $0/12^*$ & FAIL & L1 invariants \\
    advertise commit$=$lastIndex   & ok   & --    & $10/12$ & FAIL & L2 linearizability \\
    ignore \texttt{IsPaused}       & ok   & --    & $2/12$  & FAIL & L3 upstream tests \\
    \texttt{prevIndex} off-by-one  & ok   & --    & $0/12$  & FAIL & L3 upstream tests \\
    \bottomrule
  \end{tabular}
\end{table}
\noindent{\small $^\dagger$~Liveness/throughput collapse surfacing as a
sweep timeout, not a safety-invariant violation. $^*$~All $12$
linearizability seeds also timed out (reported UNKNOWN), consistent
with the throughput collapse. The control's $1/12$ linearizability
failure is the noise floor of \S\ref{sec:harness-lin}; only counts
well above it (e.g.\ commit-lie's $10/12$) are real rejections.}

\subsection{Min- vs.\ Median-Ranking on the Replicate}
\label{app:reranking}
Table~\ref{tab:reranking} re-ranks the $31$ evaluated candidates of
the replicate run (\S\ref{sec:results-loop}) under the two scoring
rules, scored against the calibration baselines (min $4905$, median
$8972$~ns/op). The rules disagree both on the winner and on how many
candidates count as improvements.

\begin{table}[h]
  \caption{Scoring-rule comparison on the replicate run's $31$
    candidates. Speedup $>1$ is faster than baseline.}
  \label{tab:reranking}
  \centering
  \small
  \begin{tabular}{llrr}
    \toprule
    Selected by & Candidate & $\min$-speedup & median-speedup \\
    \midrule
    $\min$-of-$N$ (in-loop) & \texttt{90f41d5d} & $1.06$ & $0.98$ \\
    median-of-$N$           & \texttt{5866c047} & $0.92$ & $1.60$ \\
    \midrule
    \multicolumn{2}{l}{improvements ($\ge 1.0$) under $\min$-of-$N$} & \multicolumn{2}{r}{$3/31$} \\
    \multicolumn{2}{l}{improvements ($\ge 1.0$) under median-of-$N$} & \multicolumn{2}{r}{$11/31$} \\
    \bottomrule
  \end{tabular}
\end{table}

\subsection{Candidate-Generator System Prompt}
\label{app:system-prompt}
The following is the verbatim \texttt{system\_message} supplied to the
candidate generator (the contracts of \S\ref{sec:loop} are its
\emph{hard constraints}). It is reproduced unedited except that the
specific model identity is withheld for anonymous review.

\begin{verbatim}
You are an expert in distributed systems and the etcd-io/raft Go library.
Your task is to improve a single function, (r *raft) maybeSendAppend, which
controls per-follower flow control, batching, and pipelining of MsgApp
messages on the Raft leader. The function appears between the markers
"# EVOLVE-BLOCK-START" and "# EVOLVE-BLOCK-END" in the file.

Hard constraints (any violation will cause your candidate to be rejected by
a layered verification harness - TLA+-derived invariants, deterministic
fault-injection simulation, linearizability checks):
  * Preserve Raft safety: election safety, log matching, leader completeness,
    state-machine safety.
  * Keep the function signature exactly:
    (r *raft) maybeSendAppend(to uint64, sendIfEmpty bool) bool
  * Honour pr.IsPaused(); never send when paused.
  * Always call pr.SentEntries(...) and pr.SentCommit(...) after a
    successful r.send(MsgApp).
  * Never send entries past pr.Next or beyond r.raftLog (use
    r.raftLog.entries(pr.Next, r.maxMsgSize) or a slice of its result).
  * Fall back to r.maybeSendSnapshot(to, pr) when r.raftLog.term(prevIndex)
    returns an error.

Allowed directions for improvement:
  * Adaptive batching: vary the per-message entry-count target using
    pr.Match, pr.Next, r.raftLog.committed, or r.maxMsgSize.
  * Smarter use of pr.Inflights (e.g. early returns when Inflights.Full()
    but commit index advanced).
  * Coalesce or suppress redundant empty MsgApp.
  * Pre-filter on r.raftLog.committed when nothing has actually advanced
    for this follower.
  * Reduce allocations or pointer chasing inside the hot path.

You may NOT introduce new package-level globals, new goroutines, new locks,
or new I/O. The function must remain pure with respect to its inputs.

Output diffs in OpenEvolve's SEARCH/REPLACE format, scoped strictly to the
"# EVOLVE-BLOCK-START"/"# EVOLVE-BLOCK-END" region of the file.
\end{verbatim}

%%%%%%%%%%%%%%%%%%%%%%%%%%%%%%%%%%%%%%%%%%%%%%%%%%%%%%%%%%%%

\newpage

\section*{Conference For AI Scientists 2026 - AI Involvement Checklist}
\label{checklist_ai_involvement}

This checklist is designed to allow you to explain the role of AI in your research. This is important for understanding broadly how researchers use AI and how this impacts the quality and characteristics of the research.

The checklist has two parts: \textbf{Research Stage Assessment} (items 1-4), where you rate the level of AI involvement and iteration effort at each stage of the research process, and \textbf{AI System Documentation} (items 5-6), where you provide free-text descriptions of the AI systems used and their observed limitations.

\subsection*{Research Stage Assessment}

For items 1-4, give a score from the scale below that defines the role of AI in each part of the scientific process. The scores are as follows:

\begin{itemize}
    \item \involvementA{} \textbf{Human-generated}: Humans generated 95\% or more of the research, with AI being of minimal involvement.
    \item \involvementB{} \textbf{Mostly human, assisted by AI}: The research was a collaboration between humans and AI models, but humans produced the majority (>50\%) of the research.
    \item \involvementC{} \textbf{Mostly AI, assisted by human}: The research task was a collaboration between humans and AI models, but AI produced the majority (>50\%) of the research.
    \item \involvementD{} \textbf{AI-generated}: AI performed over 95\% of the research. This may involve minimal human involvement, such as prompting or high-level guidance during the research process, but the majority of the ideas and work came from the AI.
\end{itemize}

For each research stage where AI was involved, i.e. where you selected \involvementB{}, \involvementC{}, or \involvementD{}, please also indicate the approximate level of iteration effort required using the provided iteration macros. Count substantive attempts, prompts, runs, agent trajectories, or course corrections; do not count minor wording edits.

\begin{itemize}
    \item \iterationLow{}: approximately 1-10 substantive attempts.
    \item \iterationMedium{}: approximately 10-100 substantive attempts, with some failed attempts or course corrections.
    \item \iterationHigh{}: approximately 100+ substantive attempts, with substantial exploration, failures, or trial-and-error.
    \item \iterationNA{}: not applicable because AI involvement was \involvementA{}.
    \item \iterationUnclear{}: the authors cannot reasonably estimate.
\end{itemize}

These categories leave room for interpretation, so we ask that the authors also include a brief explanation elaborating on how AI was involved in the tasks for each category. Please keep your explanation to less than 150 words.

\begin{enumerate}

    \item \textbf{Hypothesis development}: Hypothesis development includes the process by which you came to explore this research topic and research question. This can involve the background research performed by either researchers or by AI. This can also involve whether the idea was proposed by researchers or by AI.

    Answer: \involvementA{}

    Iteration Effort: \iterationNA{}

    Explanation: The research question originated from the authors'
    day-to-day work with production systems, where AI-written code was
    bottlenecked at review because no verification layer existed, and
    using a second LLM to check LLM-generated code seemed self-defeating.
    This motivated a deterministic harness that could serve as a source
    of truth for iteratively improving existing systems. We chose
    \texttt{etcd-io/raft} because it already supplies one piece of that
    harness (a published TLA$^{+}$ specification), and built outward
    from there. The idea and framing were human-generated; AI was not
    used to propose the hypothesis.

    \item \textbf{Experimental design and implementation}: This category includes design of experiments that are used to test the hypotheses, coding and implementation of computational methods, and the execution of these experiments.

    Answer: \involvementC{}

    Iteration Effort: \iterationHigh{}

    Explanation: The verification strategy was designed by the authors:
    after working with proof-based verification (e.g.\ Lean, TLA$^{+}$)
    we judged formal specifications usable as a verification layer for
    AI-generated code, then added deterministic simulation testing
    (DST), which has a strong industry track record, and a Jepsen-style
    linearizability layer. Implementation was largely AI: the agentic
    coding assistant built the harness, the deterministic simulator, the
    evaluator, and the regression tooling under human direction and
    review, and a frontier LLM was the inner candidate generator inside
    the loop. Humans set the design, the search space, and the gate; the
    AI produced most of the code and all candidate edits.

    \item \textbf{Analysis of data and interpretation of results}: This category encompasses any process to organize and process data for the experiments in the paper. It also includes interpretations of the results of the study.

    Answer: \involvementB{}

    Iteration Effort: \iterationHigh{}

    Explanation: Data processing used AI-built tooling (the evaluator's
    distribution logging, the min- vs.\ median re-ranking script, and the
    planted-bug ablation harness), but interpretation was human-led and
    corrective. The agentic assistant initially summarized results with
    optimistic single-point statistics; the authors required paired
    benchmark distributions, explicit noise-floor characterization, and a
    min-of-$N$ vs.\ median-of-$N$ comparison, which materially changed the
    reported numbers (an apparent min-of-five speedup collapsed under
    paired-median measurement). The central interpretive conclusions, that
    the in-loop metric misranks candidates and that the harness rather
    than the model is load-bearing, were drawn by the authors from
    AI-generated tables.

    \item \textbf{Writing}: This includes any processes for compiling results, methods, etc. into the final paper form. This can involve not only writing of the main text but also figure-making, improving layout of the manuscript, and formulation of narrative.

    Answer: \involvementC{}

    Iteration Effort: \iterationHigh{}

    Explanation: The manuscript was drafted with the agentic coding
    assistant, which produced initial prose for most sections, the
    LaTeX tables, and figure/layout scaffolding. The authors set the
    narrative (methodology-first framing), the claims and their scope,
    and revised extensively: tightening wording, removing over-claims,
    and reconciling numbers against the experimental logs. Multiple rounds
    of human review corrected optimistic or imprecise statements before
    they reached the final text. AI produced the majority of the draft
    text; the authors own the framing, the final claims, and verification
    of every reported number.

\end{enumerate}

\subsection*{AI System Documentation}

For items 5-6, provide a free-text description.

\begin{enumerate}
    \setcounter{enumi}{4}

    \item \textbf{AI systems used}: Describe all AI systems used in this research without naming specific products or models. Use system-level descriptors only (e.g., ``a frontier large language model'', ``a protein structure prediction system'', ``a custom multi-agent framework built on open-source models''). Include relevant details such as system type, scale, and capabilities. \textbf{Specific model and product names should only be included in the camera-ready version.}

    Description: A single frontier large language model in agentic
    coding configuration, accessed through an OpenAI-compatible chat
    completion endpoint provided by a third-party inference platform.
    The same model identity was used both as (a) the agentic coding
    assistant that constructed the verification harness, the evaluator,
    the regression tooling, and this manuscript, and (b) the inner
    candidate generator inside the evolutionary optimization loop. An
    open-source evolutionary code-optimization framework derived from
    AlphaEvolve provided the MAP-Elites candidate database, parent
    sampler, and prompt-templating layer; we did not modify the
    framework. The Porcupine linearizability checker is used as the
    consistency oracle in the harness but is a deterministic
    algorithmic component, not an AI system.

    \item \textbf{Observed AI Limitations}: What specific limitations and failure modes have you observed when using AI as a partner or lead author?

    Description: Three failure modes recurred. (1)~\emph{Method
    invention}: the inner candidate generator occasionally proposed
    code referencing methods that do not exist on the relevant
    structs (e.g.\ a hallucinated \texttt{IsHealthy} on the progress
    tracker); these were caught at the compile layer of the harness
    rather than ending up in deployed code. (2)~\emph{Score-metric
    over-fitting}: when the evaluator scored on a min-of-five
    benchmark, the generator produced candidates whose ``speedup'' was
    largely the result of single-run jitter in the lower tail; the
    apparent fifteen-percent improvement reduced to a paired-median
    four percent under more rigorous measurement. (3)~\emph{Optimistic
    self-reporting}: when asked to summarise experimental outputs,
    the agentic coding assistant initially reported the
    optimistic-end statistic; subsequent prompts requiring paired
    distributions and explicit noise-floor characterisation produced
    materially different numbers. We did not measure cross-model
    stability, and the underlying inference service is stochastic;
    different runs of the same loop are not expected to produce
    identical candidates.

\end{enumerate}

\newpage

\section*{Conference For AI Scientists 2026 - Reproducibility and Responsibility Checklist}
\label{checklist_reproducibility}

\begin{enumerate}

\item {\bf Claims}
    \item[] Question: Do the main claims made in the abstract and introduction accurately reflect the paper's contributions and scope?
    \item[] Answer: \answerYes{}
    \item[] Justification: The abstract and \S\ref{sec:intro} enumerate four contributions (harness, calibration, evolutionary-search wrapper, empirical study) that map directly to \S\ref{sec:harness}, \S\ref{sec:results-loop}, \S\ref{sec:loop}, and \S\ref{sec:results}.

\item {\bf Limitations}
    \item[] Question: Does the paper discuss the limitations of the work performed by the authors?
    \item[] Answer: \answerYes{}
    \item[] Justification: \S\ref{sec:limits} discusses oracle coverage as a lower bound, the fault catalogue's omissions (crash-restart, slow disk), the synthetic-performance caveat, the harness-vs-deployment distinction, generator stochasticity, and benchmark noise floors.

\item {\bf Theory assumptions and proofs}
    \item[] Question: For each theoretical result, does the paper provide the full set of assumptions and a complete (and correct) proof?
    \item[] Answer: \answerNA{}
    \item[] Justification: The paper does not include theoretical results. The Raft invariants we check come from prior published work~\cite{ongaro2014search} and the in-tree TLA$^{+}$ specification; we do not extend or re-prove them.

    \item {\bf Experimental result reproducibility}
    \item[] Question: Does the paper fully disclose all the information needed to reproduce the main experimental results of the paper to the extent that it affects the main claims and/or conclusions of the paper (regardless of whether the code and data are provided or not)?
    \item[] Answer: \answerYes{}
    \item[] Justification: \S\ref{sec:setup} records the pinned upstream commit (\texttt{etcd-io/raft @ 67d129e8}, 2026-05-26), the Go and tooling versions (Go \texttt{1.26.3}, OpenEvolve \texttt{0.2.27}, Porcupine \texttt{v1.1.0}, OpenJDK \texttt{17}), the hardware configuration (dedicated 16-vCPU \texttt{x86\_64} cloud host, 32~GB RAM, Ubuntu \texttt{24.04}), the seeded fault parameters, the seed set, and the benchmark methodology (Go's standard \texttt{go test -bench=\dots{} -benchtime=1s -count=5} with min-reduction). The single upstream patch enabling deterministic seeded election timeouts is described in \S\ref{sec:harness}.

\item {\bf Open access to data and code}
    \item[] Question: Does the paper provide open access to the data and code, with sufficient instructions to faithfully reproduce the main experimental results, as described in supplemental material?
    \item[] Answer: \answerYes{}
    \item[] Justification: The harness source, evaluator, OpenEvolve configuration, prompt templates, and SQLite candidate database will be released as the verification artefact on acceptance, alongside the per-seed reproduction scripts referenced in \S\ref{sec:results}. The release will be anonymised at submission and de-anonymised at camera-ready per CAISc policy.

\item {\bf Experimental setting/details}
    \item[] Question: Does the paper specify all the training and test details (e.g., data splits, hyperparameters, how they were chosen, type of optimizer, etc.) necessary to understand the results?
    \item[] Answer: \answerYes{}
    \item[] Justification: \S\ref{sec:setup} and \S\ref{sec:loop} detail cluster size, tick budget, fault probabilities, scenario parameters, search space, generator temperature, max-tokens, and evaluator timeouts. The OpenEvolve config YAML is included in the artefact.

\item {\bf Experiment statistical significance}
    \item[] Question: Does the paper report error bars suitably and correctly defined or other appropriate information about the statistical significance of the experiments?
    \item[] Answer: \answerYes{}
    \item[] Justification: \S\ref{sec:results-variants} (Table~\ref{tab:bench-paired}) reports the min, median, and max of a paired ten-run benchmark distribution rather than single-point speedups, and \S\ref{sec:limits} explicitly characterises the bimodal benchmark distribution and a $\sim$17\% linearizability false-positive noise floor below which apparent improvements are not interpretable as real. We do not perform parametric hypothesis testing because the underlying distribution violates normality.

\item {\bf Experiments compute resources}
    \item[] Question: For each experiment, does the paper provide sufficient information on the computer resources (type of compute workers, memory, time of execution) needed to reproduce the experiments?
    \item[] Answer: \answerYes{}
    \item[] Justification: \S\ref{sec:setup} states the single-host commodity workstation context and the embarrassingly-parallel-by-seed nature of the harness. \S\ref{sec:results-cost} reports the per-iteration $\sim$57~s wall-clock decomposition, the $\sim$28-minute total run time, and the approximate ten-dollar inference spend.

\item {\bf Research Ethics}
    \item[] Question: Does the research conducted in the paper conform with the \href{https://neurips.cc/Conferences/2026/MainTrackHandbook}{NeurIPS Code of Ethics}, except where CAISc's AI authorship policies apply?
    \item[] Answer: \answerYes{}
    \item[] Justification: The research uses no human subjects, no personally identifiable data, and no proprietary corpora. The AI Involvement Checklist documents the role of the language model. The work targets a published open-source library; we do not introduce undisclosed modifications and we propose no production deployment.

\item {\bf Broader impacts}
    \item[] Question: Does the paper discuss both potential positive societal impacts and negative societal impacts of the work performed?
    \item[] Answer: \answerYes{}
    \item[] Justification: Positive: harness-gated agentic loops lower the cost of safety-preserving optimisation in safety-critical infrastructure, and the verification-pyramid pattern generalises to other distributed-systems libraries with published TLA$^{+}$ specifications. Negative: the same loop could be misused to produce subtly-incorrect code that nonetheless clears a weak harness; we mitigate by emphasising in \S\ref{sec:limits} that harness-acceptance is not deployment-readiness and that the oracle's coverage is a lower bound on bug-finding capability.

\end{enumerate}


\end{document}
